\documentclass[11pt]{article}
\usepackage[margin=1in]{geometry}
\usepackage{hyperref}
\usepackage{enumitem}
\usepackage{booktabs}
\usepackage{xcolor}
\title{ProbOS --- Month 3, Week 13\\PDSL v0.2: Control Flow}
\author{Nisong Monyimba}
\date{\today}

\begin{document}
\maketitle

\section*{Week 13 goal}
Extend PDSL from v0.1 to v0.2 by adding comparison operators and
conditional expressions -- the first control-flow constructs PDSL
has ever supported.

\textbf{Scope confirmed via direct inspection of \texttt{grammar.lark}
before writing this plan:} PDSL v0.1 has ZERO comparison operators
and ZERO conditional constructs -- only arithmetic (+, -, *, /, **,
unary +/-) and function calls.

\textbf{Design decision, stated up front:} this adds
EXPRESSION-level conditionals (\texttt{if cond then expr else expr}),
NOT statement-level control flow or loops. PDSL describes a model's
per-step drift declaratively; the outer simulation engine already
handles the step-by-step loop. An expression-level conditional (for
clamping, thresholding, or piecewise physical behavior) is the right
granularity.

\section*{Day-by-day plan}

\subsection*{Day 1 -- Grammar design}
\begin{itemize}[leftmargin=*]
  \item Add comparison operators to \texttt{grammar.lark}: \texttt{<},
    \texttt{>}, \texttt{<=}, \texttt{>=}, \texttt{==}, \texttt{!=} --
    at a new precedence level between \texttt{expr} and a new
    top-level \texttt{condition} rule
  \item Add a conditional expression rule:
    \texttt{"if" condition "then" expr "else" expr -> conditional}
  \item Test the grammar in isolation (Lark parse-tree inspection)
    before touching \texttt{parser.py} or \texttt{codegen.py}
\end{itemize}

\subsection*{Day 2 -- AST nodes and parser}
\begin{itemize}[leftmargin=*]
  \item Add \texttt{Comparison} and \texttt{Conditional} AST node
    types to \texttt{ast\_nodes.py}, matching existing node style
  \item Extend \texttt{parser.py} to build these new AST nodes
\end{itemize}

\subsection*{Day 3 -- Code generation}
\begin{itemize}[leftmargin=*]
  \item Extend \texttt{codegen.py} to emit correct, VECTORISED
    Python: a naive \texttt{if/else} would not vectorise correctly
    over N particles -- must generate \texttt{xp.where(cond,
    expr\_true, expr\_false)}, matching the vectorisation discipline
    already used throughout the hand-written kernel models
\end{itemize}

\subsection*{Day 4 -- Test coverage}
\begin{itemize}[leftmargin=*]
  \item Grammar-level, codegen-level, and end-to-end tests: a
    genuinely new PDSL model using a conditional compiles and runs
    correctly through the full pipeline, matching a hand-written
    Python equivalent
\end{itemize}

\subsection*{Day 5 -- Existing model regression check}
\begin{itemize}[leftmargin=*]
  \item Confirm ALL existing PDSL v0.1 models/examples still parse
    and compile correctly -- strictly additive, not breaking
\end{itemize}

\subsection*{Day 6 -- pre\_week\_audit.sh + CI}
\begin{itemize}[leftmargin=*]
  \item Full audit pass
\end{itemize}

\subsection*{Day 7 -- Documentation + retrospective}
\begin{itemize}[leftmargin=*]
  \item Update \texttt{docs/vision/main.tex} if it references PDSL
    version status
  \item Retrospective appended once complete
\end{itemize}

\section*{Exit criteria}
PDSL v0.2 correctly parses, compiles, and executes conditional
expressions with comparison operators, generating correctly-
vectorised NumPy/CuPy code (\texttt{xp.where}), with zero regressions
to any existing v0.1 model.

\end{document}
